%------------------------------------------------------------------------------%
\begin{question}
  Determine a posição relativa entre a reta e o plano \\
  e o ponto de interseção, se ele existir
  \[
    r\colon
    \begin{cases}
      x = 1 +  t \\
      y = 2 + 2t \\
      z = 1 -  t
    \end{cases}
  \]
  \[
    \gamma\colon x + y - z + 10 = 0
  \]
\end{question}

%------------------------------------------------------------------------------%
\begin{resolution}{Posição Relativa}
  \onslide<+->{O vetor diretor da reta é \qquad
  \(
    \vec{v} = (1, 2, -1)
  \)}

  \onslide<+->{O vetor normal do plano é \quad
  \(
    \vec{n} = (1, 1, -1)
  \)}
  \begin{align*}
    \onslide<+->{\vec{v} \cdot \vec{n}}
    \onslide<+->{& = 1 \times 1 + 2\times 1 + (-1)(-1)} \\
    \onslide<+->{& = 1 + 2 - 1} \\
    \onslide<+->{& = 2} \\
    \onslide<+->{& \neq 0}
  \end{align*}
  \onslide<+->{a reta é secante ao plano}
\end{resolution}

%------------------------------------------------------------------------------%
\begin{resolution}{Ponto de Interseção}
\begin{columns}
\column{0.35\textwidth}
  \onslide<+->{Substituir a equação da reta
  \[
    \begin{cases}
      x = 1 +  t \\
      y = 2 + 2t \\
      z = 1 -  t
    \end{cases}
  \]}
  \onslide<+->{na equação do plano}
\column{0.5\textwidth}
  \begin{align*}
    \onslide<+->{x + y - z + 10              & = 0 \\}
    \onslide<+->{(1+t) + (2+2t) - (1-t) + 10 & = 0 \\}
    \onslide<+->{t + 2t + t + 1 + 2 - 1 + 10 & = 0 \\}
    \onslide<+->{4t + 12                     & = 0 \\}
    \onslide<+->{t                           & = \frac{12}{4} \\}
    \onslide<+->{                            & = 3}
  \end{align*}
\end{columns}
\end{resolution}

%------------------------------------------------------------------------------%
\begin{resolution}{Ponto de Interseção}
  \onslide<+->{Coordenadas do ponto com $t=0$}
  \begin{align*}
    \onslide<+->{x & = 1 + t}
    \onslide<+->{  \;=\; 1 + 3}
    \onslide<+->{  \;=\; 4 \\}
    \onslide<+->{y & = 2 + 2t}
    \onslide<+->{  \;=\; 2 + 2\times 3}
    \onslide<+->{  \;=\; 8 \\}
    \onslide<+->{z & = 1 - t}
    \onslide<+->{  \;=\; 1 - 3}
    \onslide<+->{  \;=\; -2}
  \end{align*}

  \onslide<+->{\[
    P = (4, 8, -2)
  \]}
\end{resolution}

%------------------------------------------------------------------------------%
\endinput

